\section{\emph{Prompt} e registro da Etapa II}\label{apx:prompt}

O enunciado pede que o \emph{prompt} e a versão do modelo sejam documentados~\cite{enunciado}. Esta seção fixa o
\emph{prompt} antes da execução. O texto é em inglês e usa apenas caracteres ASCII, para evitar problemas de
codificação em APIs; os campos entre chaves são preenchidos pelo código.

\subsection{Mensagens}
\begin{lstlisting}
=== SYSTEM MESSAGE ===
You are an expert in server monitoring and anomaly detection.

You will receive one observation of a server: 38 anonymous
resource-utilization and activity metrics, named m1 to m38. Every value is
scaled to the range [0, 1]. Metrics are sampled once per minute. Their names
and units are not available.

Task: decide whether the observation belongs to an anomalous period (a server
incident) or to normal operation.

Rules:
- Use only the information given in this conversation.
- Anomalies are rare: roughly 4 out of every 100 observations.
- Reply with one JSON object and nothing else:
  {"label": 0 or 1, "score": a number between 0 and 1}
  label = 1 if the observation is anomalous, 0 if it is normal.
  score = your confidence that the observation is anomalous
          (0 = certainly normal, 1 = certainly anomalous).

=== USER MESSAGE, zero-shot ===
Observation:
{observation}

=== USER MESSAGE, few-shot (k = 3, 5, 10 or 20 labeled examples) ===
Labeled examples from the same kind of server (k = {k}):

Example 1:
{observation_1}
Answer: {"label": 0, "score": 0.0}

Example 2:
{observation_2}
Answer: {"label": 1, "score": 1.0}

... (up to Example {k}, in random order)

Now classify this observation:
{observation}

=== OPTIONAL BLOCK, reference profile (zero-shot variant) ===
Typical values of each metric during normal operation of this server
(median / 95th percentile):
m1: {median_1} / {p95_1}
...
m38: {median_38} / {p95_38}
\end{lstlisting}

\subsection{Serialização das observações}
Cada observação é escrita em uma linha, com os 38 valores no formato \texttt{mj=v}, separados por espaço e com
4 casas decimais. O exemplo abaixo é a linha 15.000 do arquivo de teste da máquina 1-1 (rótulo 0), quebrada em cinco linhas apenas para caber na página; no \emph{prompt}, é uma linha só. A linha 16.100 da mesma máquina
(rótulo 1) começa com \texttt{m1=0.5161}, acima do máximo de treino do atributo (0,4946), e serve de exemplo de
observação anômala (Seção~\ref{sec:exemplos}). Ambas são ilustrações da \emph{forma} do texto, não exemplos do
protocolo.
\begin{lstlisting}
m1=0.0968 m2=0.0794 m3=0.0702 m4=0.0627 m5=0.0000 m6=0.8974 m7=0.2750 m8=0.0000
m9=0.0288 m10=0.0000 m11=0.1066 m12=0.1081 m13=0.0274 m14=0.0575 m15=0.1250 m16=0.1246
m17=0.0000 m18=0.0000 m19=0.1483 m20=0.1009 m21=0.1254 m22=0.0984 m23=0.8330 m24=0.0126
m25=0.0493 m26=0.0117 m27=0.0000 m28=0.1029 m29=0.0000 m30=0.0201 m31=0.0897 m32=0.0839
m33=0.0004 m34=0.0001 m35=0.0810 m36=0.0811 m37=0.0000 m38=0.0000
\end{lstlisting}

\subsection{Regras de execução}
\begin{itemize}
  \item \textbf{Exemplos:} para cada par $(k,\text{semente})$ sorteia-se de $L_{\mathrm{tr}}$ uma lista fixa de $k$
        exemplos, com $k=3$: 2~normais e 1~anômalo; $k=5$: 3+2; $k=10$: 5+5; $k=20$: 10+10, em ordem
        aleatória. A mesma lista é usada para todas as observações da execução, o que permite o uso de
        \emph{prompt caching}, se a API o oferecer.
  \item \textbf{Parâmetros:} temperatura~0 (ou a menor disponível) e limite curto de \emph{tokens} de saída; usar o
        modo de saída JSON quando existir.
  \item \textbf{Respostas inválidas:} se a resposta não for um JSON válido, repete-se a chamada uma vez; persistindo
        o erro, a observação conta como rótulo~0 e escore~0, e a taxa de respostas inválidas é reportada.
  \item \textbf{Escore:} o escore autodeclarado não é calibrado, de modo que a AUC-ROC calculada sobre ele é
        apenas indicativa. Se a API fornecer \emph{logprobs}, a probabilidade do \emph{token} do rótulo é uma
        alternativa.
  \item \textbf{Número de chamadas:} para $n$ observações por máquina, são $3\times[1+4\times3]\times n=39\,n$
        chamadas (3 máquinas, 1 configuração \emph{zero-shot} e 4 valores de $k$ com 3 sementes cada); por
        exemplo, 19.500 chamadas com $n=500$. O custo total segue de multiplicar esse número pelos \emph{tokens}
        médios por chamada e pelos preços da Tabela~\ref{tab:registro-llm}.
  \item \textbf{Contaminação:} o SMD é público desde 2019~\cite{su2019,omnianomaly}, e não é possível excluir
        que faça parte dos dados de pré-treinamento do modelo; isso deve ser mencionado na discussão.
\end{itemize}

\subsection{Registro da execução}
\begin{table}[htbp]
\centering
\caption{Registro da Etapa II (preencher durante a execução).}
\label{tab:registro-llm}
\small
\begin{tabularx}{\textwidth}{@{}lL@{}}
\toprule
Campo & Valor\\
\midrule
Provedor e modelo & \apreencher\\
Identificador exato da versão (cadeia da API) & \apreencher\\
Data(s) de execução & \apreencher\\
Temperatura e limite de \emph{tokens} de saída & 0 e valor curto (proposta); valores usados: \apreencher\\
Observações avaliadas por máquina ($n$) & 500 a 2.000 (proposta); valor usado: \apreencher\\
\emph{Tokens} médios por chamada (entrada / saída) & \apreencher\\
Preço por milhão de \emph{tokens} (entrada / saída) & \apreencher\\
Tempo por 100 amostras & \apreencher\\
Custo estimado por 100 amostras & \apreencher\\
Respostas inválidas (\%) & \apreencher\\
\bottomrule
\end{tabularx}
\end{table}
